\section{Gestión de la Información (Cumplimiento ISO 26515)}
\label{sec:s5_gestion}

\subsection{Publicación de las Notas de la Versión (Release Notes)}
\label{ssec:s5_gi_releasenotes}
Se publicaron las Notas de la Versión integrales del incremento 2.0.0, consolidando las
funcionalidades entregadas a lo largo de los seis sprints, los cambios críticos
(\textit{breaking changes}) y las acciones de seguridad recomendadas.

\subsection{Empaquetado y Versionado de la Base de Conocimiento}
\label{ssec:s5_gi_empaquetado}
La base de conocimiento (diagramas, arquitecturas y guías de uso) se empaquetó y versionó
directamente embebida dentro del repositorio oficial del código fuente, cerrando el ciclo de
\textit{Documentation as Code}. La presente Memoria P.D.S. se etiqueta como
\comando{PDS-ETHOSHUB-2.0.0}.

\begin{TablaHSDoc}
    EPIC-06 & Visibilidad y SEO & Sí (S5) & Sí (S5) \\ \hline
    Todas & Cierre de versión 2.0.0 & Sí (S5) & Sí (S5) \\ \hline
\end{TablaHSDoc}

\subsection{Productos de Información del Cierre (ISO 15289)}
\label{ssec:s5_gi_productos}
La tabla~\ref{tab:s5_productos} relaciona los productos de información del cierre de versión con
su artefacto de origen.

\begin{table}[!ht]
    \centering
    \renewcommand{\arraystretch}{1.3}
    \fontsize{9}{10.8}\selectfont\sffamily
    \begin{tabularx}{\textwidth}{|X|X|}
        \hline
        \rowcolor{colorPrimario}
        \textcolor{white}{\textbf{Producto de información}} & \textcolor{white}{\textbf{Artefacto de origen}} \\ \hline
        Release Notes 2.0.0 & Historial de incrementos (Sprints 0--5). \\ \hline
        \rowcolor{grisTecnico}
        Guía de despliegue monolítico & \comando{build.sh}, \comando{Dockerfile}, \comando{docker-compose.yml}. \\ \hline
        Base de conocimiento empaquetada & Diagramas, arquitecturas y guías embebidas en el repo. \\ \hline
    \end{tabularx}
    \caption{Productos de información del Sprint 5 y su trazabilidad.}
    \label{tab:s5_productos}
\end{table}

\clearpage
